\Needspace{8\baselineskip}\originalchapter{2003 作业一: 分布与电离}
题源: I. H. Hutchinson, 10 Sep 2003, due 17 Sep. 本章及随后2003六套作业均无当期官方答案文件, 以下为编者独立补解, 不称MIT解答.
\begin{exercise}\label{ex:03-1-1}
原题1. 平均量定义 $\langle G\rangle=\int Gf\dd^3v/\int f\dd^3v$, 三维 Maxwell $f=n[m/(2\pi T)]^{3/2}\exp[-mv^2/(2T)]$. (a)证明归一为n; (b)求 $\langle v_x\rangle$; (c)求 $\langle v^2\rangle$ 和平均动能; (d)求平均速率. 各向同性使三 Cartesian 方向等价.
\end{exercise}
\begin{solution}
用 $\int_{-\infty}^{\infty}\e^{-av^2}\dd v=\sqrt{\pi/a}$, 三积分相乘抵消前因子, 得n. 奇函数给 $\langle v_x\rangle=0$. 对 Gaussian 积分对a求导得单轴 $\langle v_x^2\rangle=T/m$, 因而 $\langle v^2\rangle=3T/m$, $\langle mv^2/2\rangle=3T/2$. 球坐标速率平均含 $4\pi\int_0^\infty v^3f(v)\dd v/n$, 代 $u=mv^2/(2T)$ 得 $\sqrt{8T/(\pi m)}$. 平均矢量速度零不等于平均速率零.
\end{solution}
\begin{exercise}\label{ex:03-1-2}
原题2. 对任意仅依 $|v|$ 的各向同性分布, 只取某轴速度分量为正的粒子, 证明单侧撞击面通量为 $n\langle|v|\rangle/4$.
\end{exercise}
\begin{solution}
法向取z, $\Gamma=2\pi\int_0^\infty v^3f(v)\dd v\int_0^{\pi/2}\sin\theta\cos\theta\dd\theta=\pi\int v^3f\dd v$. 全空间的 $n\langle|v|\rangle=4\pi\int v^3f\dd v$, 得1/4. 这个结果需要各向同性, 有向束流不能套用.
\end{solution}
\begin{exercise}\label{ex:03-1-3}
原题3. 查出并记住下列量, 除注明外以SI两位有效数字: c、\ensuremath{\epsilon}0、\ensuremath{\mu}0、$1/\sqrt{\epsilon_0\mu_0}$、电子电荷大小e、电子与质子质量、1eV对应K、呼吸空气粒子密度与torr/Pa压强、水的粒子密度、氢电离能、Gauss/Tesla和cm$^{-3}$/m$^{-3}$换算.
\end{exercise}
\begin{solution}
$c=3.0\times10^8$ m/s; $\epsilon_0=8.9\times10^{-12}$ F/m; $\mu_0\simeq1.3\times10^{-6}$ H/m; $1/\sqrt{\epsilon_0\mu_0}=c$. $e=1.6\times10^{-19}$ C, $m_e=9.1\times10^{-31}$ kg, $m_p=1.7\times10^{-27}$ kg; $1$ eV=$1.2\times10^4$ K. 取300K/1atm空气, $n=P/(k_BT_K)=2.4\times10^{25}$ m$^{-3}$, P=760torr=$1.0\times10^5$ Pa. 常温水按分子数为 $3.3\times10^{28}$ m$^{-3}$, 若按原子数为三倍, 须说明粒子定义. 氢基态电离能13.6eV; $1$G=$10^{-4}$T; $1$cm$^{-3}=10^6$m$^{-3}$. 新SI中\ensuremath{\mu}0不再定义为精确 $4\pi10^{-7}$, 此精度取近似无影响.
\end{solution}
\begin{exercise}\label{ex:03-1-4}
原题4. Saha方程 $n_en_i/n_0=2(2\pi m_eT)^{3/2}h^{-3}\e^{-\chi_i/T}$, Debye长度远小于放电尺度. 当\ensuremath{\chi}i=13eV、总压强1atm, 求气体50\%电离时约T, 以及电子密度与同压室温气体密度之比.
\end{exercise}
\begin{solution}
50\%按原子核计: $n_i=n_e=n_0$, 总压强 $P=(n_0+n_i+n_e)T=3n_eT$. 代入Saha得到一个单调根方程 $P/(3T)=2(2\pi m_eT)^{3/2}h^{-3}\e^{-13\,\mathrm{eV}/T}$. 实算 $T=1.2104$eV=$1.4047\times10^4$K. 密度比 $n_e/[P/(k_B300\mathrm K)]=300k_B/(3T)=.007119$. 总压强计入电子, 否则会错用2而非3. Saha需要局部热平衡; 稀薄非平衡放电的50\%不必由此T给出.
\end{solution}
